\documentclass[a4paper]{arrowhead}

\usepackage[yyyymmdd]{datetime}
\usepackage{etoolbox}
\usepackage[T1]{fontenc}
\usepackage[utf8]{inputenc}
\usepackage{multirow}
\usepackage{longtable}
\usepackage{amssymb}
\usepackage{float}
\usepackage{tikz}
\usetikzlibrary{positioning,arrows.meta,fit,calc,shapes.misc,shapes.geometric}

\renewcommand{\dateseparator}{-}
\setlength{\emergencystretch}{3em}
\lstset{basicstyle=\ttfamily\footnotesize,breaklines=true,columns=fullflexible,frame=single,rulecolor=\color{gray!50},backgroundcolor=\color{gray!6}}
\tikzset{
  ahsys/.style={draw=ArrowheadBlue,very thick,fill=ArrowheadBlue!6,rounded corners=2pt,minimum width=3.2cm,minimum height=1.1cm,align=center,font=\small\bfseries},
  ahext/.style={draw=ArrowheadGrey,thick,dashed,fill=gray!5,rounded corners=2pt,minimum width=3.2cm,minimum height=1.0cm,align=center,font=\small},
  ahsrv/.style={draw=ArrowheadPurple,thick,fill=ArrowheadPurple!8,rounded corners=6pt,minimum height=0.55cm,inner xsep=4pt,align=center,font=\scriptsize\ttfamily},
  ahcsm/.style={draw=ArrowheadTeal,thick,fill=ArrowheadTeal!8,rounded corners=6pt,minimum height=0.55cm,inner xsep=4pt,align=center,font=\scriptsize\ttfamily},
  ahdb/.style={draw=ArrowheadOrange,thick,fill=ArrowheadOrange!8,cylinder,shape border rotate=90,aspect=0.25,minimum height=0.9cm,align=center,font=\scriptsize},
  ahflow/.style={-{Stealth[length=2.2mm]},thick,ArrowheadGrey},
  ahlabel/.style={font=\scriptsize,fill=white,inner sep=1pt,align=center},
}

%% Special references
\newcommand{\mref}[1]{{\textcolor{ArrowheadPurple}{\hyperref[sec:model:#1]{#1}}}}
\newcommand{\oref}[1]{{\textcolor{ArrowheadBlue}{\hyperref[sec:operations:#1]{#1}}}}
\newcommand{\pdef}[1]{{\textcolor{ArrowheadGrey}{#1\label{sec:model:primitives:#1}\label{sec:model:primitives:#1s}\label{sec:model:primitives:#1es}}}}
\newcommand{\pref}[1]{{\textcolor{ArrowheadGrey}{\hyperref[sec:model:primitives:#1]{#1}}}}

\newrobustcmd\osubsection[3]{
  \addtocounter{subsection}{1}
  \addcontentsline{toc}{subsection}{\protect\numberline{\thesubsection}operation \textcolor{ArrowheadBlue}{#1}}
  \subsection*{
    \thesubsection\quad
    operation
    \textcolor{ArrowheadBlue}{#1}
    (\notblank{#2}{\mref{#2}}{})
    \notblank{#3}{: \mref{#3}}{}
  }
  \label{sec:operations:#1}
}
\newrobustcmd\msubsection[2]{
  \addtocounter{subsection}{1}
  \addcontentsline{toc}{subsection}{\protect\numberline{\thesubsection}#1 \textcolor{ArrowheadPurple}{#2}}
  \subsection*{\thesubsection\quad#1 \textcolor{ArrowheadPurple}{#2}}
  \label{sec:model:#2} \label{sec:model:#2s} \label{sec:model:#2es}
}

\begin{document}

%% Arrowhead Document Properties
\ArrowheadTitle{workflowInstance}
\ArrowheadServiceID{workflowInstance}
\ArrowheadType{Service Description}
\ArrowheadTypeShort{SD}
\ArrowheadVersion{5.4.0}
\ArrowheadDate{2026-11-23}
\ArrowheadAuthor{Documentation maintenance}
\ArrowheadStatus{REVIEW DRAFT}
\ArrowheadContact{Repository maintainers}
\ArrowheadFooter{\href{www.arrowhead.eu}{www.arrowhead.eu}}
\ArrowheadSetup
%%

%% Front Page
\begin{center}
  \vspace*{1cm}
  \huge{\arrowtitle}

  \vspace*{0.2cm}
  \LARGE{\arrowtype}
  \vspace*{1cm}

  \Large{Service ID: \textit{``\arrowid''}}
  \vspace*{\fill}

  

  \vspace*{1cm}
  \vspace*{\fill}

  % Front Page Abstract
  \begin{abstract}
    This document describes the \texttt{workflowInstance} service of the Eclipse Arrowhead workflow support systems, which starts, observes and aborts workflow sessions. It is provided by the WorkflowManager system, service version 1.0.0, documentation edition 5.4.0.
  \end{abstract}

  \vspace*{1cm}
\end{center}

\newpage
%%

%% Table of Contents
\tableofcontents
\newpage
%%

\section{Overview}
\label{sec:overview}
The \texttt{workflowInstance} service starts and observes \emph{sessions}. A session is one
execution of one immutable recipe version for one subject: a product unit, a control period or a batch. Its state is
$z = (M, M_A, \beta, \varepsilon, \xi, status)$~\cite{varga2026}: the step and action markings, the service bindings,
the executor assignment, the stored results and the status. The WorkflowManager advances the state with the transition
rules of the model, dispatches enabled steps to WorkflowExecutors, and exposes the markings and the full trace through
this service.

The rest of this document is organized as follows.
In Section~\ref{sec:operations}, we describe the abstract message operations provided by the service.
In Section~\ref{sec:model}, we end the document by presenting the data types used by the mentioned operations.

\subsection{Significant Prior Art}
The session state and its eight transition rules (A-Start, S-Enable, S-Bind, S-Done, A-Finish, Z-Close,
S-FailBind, S-FailInvoke) are those of Varga~\cite{varga2026}, which also proves that a failure-free run of a
well-formed recipe takes exactly $3|S| + 2|A|$ transitions (Theorem~1) and always progresses (Theorem~2). The
product-driven start, where the product hands its recipe to the work centre, is Derhamy's Intelligent Product
pattern~\cite{derhamy2018workflow}.

\subsection{How This Service Is Meant to Be Used}
A consumer starts a session with \oref{start}, naming a registered recipe (and optionally its version), the
session parameters $\rho$ and the dispatch mode. Alternatively it names a \texttt{product} system: the WorkflowManager
then retrieves the recipe fragment and parameters from the product's \texttt{workflowProduct} service and stores the
results back when the session completes. The owner follows the session with \oref{get-status} and \oref{get-trace} and may
stop it with \oref{abort}. \oref{query} lists sessions, e.g. all \texttt{running} ones.

\subsection{Important Delimitations and Dependencies}
\label{sec:delimitations}
\begin{itemize}
\item Failure is absorbing: a failed session never resumes. A new session must be started.
\item Theorems~1--2 hold only for sessions without extensions. The \texttt{extensions} field and \texttt{withinModel} in the trace say whether a session used retry/rebind, abort, timeout or guarded skipping.
\item Only the owner (the starting system) or a sysop may abort a session.
\item Sessions run in one local cloud. Delegation of actions to WorkflowManagers in other clouds is not implemented.
\item The \texttt{event} dispatch mode is accepted but handled as \texttt{pull}.
\end{itemize}

\paragraph{Authorization.} Each operation is authorized separately. \texttt{abort} additionally requires the caller to be the session owner or a sysop.

\newpage
\section{Service Interface}
\label{sec:operations}
This section describes the interface of the \texttt{workflowInstance} service. Each subsection names an abstract operation,
an input type and an output type, in that order. The input type is named inside parentheses, while the output type is
preceded by a colon. Input and output types are only denoted when accepted or returned, respectively, by the operation
in question. All abstract data types named in this section are defined in Section~\ref{sec:model}. The concrete
interface (HTTP/TLS/JSON) is specified in the IDD of the same name.

\begin{figure}[ht!]
  \centering
  \resizebox{0.92\textwidth}{!}{
\begin{tikzpicture}
\node[ahsys] (prov) at (4.6,0) {WorkflowManager\\[-1pt]{\normalfont\scriptsize $\ll$provider$\gg$}};
\node[ahsrv] (srv) at (1.6,0) {workflowInstance};
\draw[thick,ArrowheadPurple] (srv.east) -- (prov.west);
\node[ahext] (c0) at (-3.0,1.20) {Intelligent Product\\[-1pt]{\scriptsize $\ll$consumer$\gg$}};
\draw[ahflow] (c0.east) -- (srv.west);
\node[ahext] (c1) at (-3.0,0.00) {Line controller\\[-1pt]{\scriptsize $\ll$consumer$\gg$}};
\draw[ahflow] (c1.east) -- (srv.west);
\node[ahext] (c2) at (-3.0,-1.20) {Upper WorkflowManager\\[-1pt]{\scriptsize $\ll$consumer$\gg$}};
\draw[ahflow] (c2.east) -- (srv.west);
\end{tikzpicture}
  }
  \caption{SysML block view of the \texttt{workflowInstance} service: provider and typical consumers.}
  \label{fig:SD}
\end{figure}

The following interface operations are available.

{\small
\begin{longtable}{|p{4.4cm}|p{5.6cm}|p{5.8cm}|}

\hline
\rowcolor{gray!33} Operation & Input & Output \\ \hline
\endhead
\oref{start} & \mref{StartSessionRequest} & \mref{Session} \\ \hline
\oref{abort} & \mref{SessionRef} & \mref{Session} \\ \hline
\oref{query} & \mref{SessionFilter} & \mref{SessionList} \\ \hline
\oref{get-status} & \mref{SessionRef} & \mref{Session} \\ \hline
\oref{get-trace} & \mref{SessionRef} & \mref{Trace} \\ \hline
\end{longtable}
}

\osubsection{start}{StartSessionRequest}{Session}
Creates a session of an active recipe version (or of the product's recipe), fires the first transitions and dispatches the enabled steps.
\osubsection{abort}{SessionRef}{Session}
Aborts a running session (extension x-abort). Its active work items become \texttt{ABORTED} and their executors receive \texttt{workflowExecution.abort}; leases are released when \texttt{workflow.lease.release-on-failure} is set.
\osubsection{query}{SessionFilter}{SessionList}
Lists sessions, newest first, optionally in one status.
\osubsection{get-status}{SessionRef}{Session}
Returns the markings and status of one session.
\osubsection{get-trace}{SessionRef}{Trace}
Returns the run: every transition, $\xi$, $\beta$, $\varepsilon$, and whether the run stayed within the proven model.

\newpage
\section{Information Model}
\label{sec:model}
Here, all data objects that can be part of the \texttt{workflowInstance} service are listed in alphabetic order. Each
subsection, which describes one type of object, begins with the \textit{struct} keyword (a collection of named fields,
each with its own data type), the \textit{enum} keyword (a closed set of values) or the \textit{alias} keyword. As a
complement to the explicitly defined types in this section, there is also a list of implicit primitive types in
Section~\ref{sec:model:primitives}.

\msubsection{enum}{DispatchMode}
How a step reaches an executor.

{\small
\begin{longtable}{|p{3.9cm}|p{11.9cm}|}

\hline
\rowcolor{gray!33} Value & Meaning \\ \hline
\endhead
\texttt{push} & The WorkflowManager selects an executor and calls \texttt{workflowExecution.start} (default). \\ \hline
\texttt{pull} & The WorkflowManager publishes the work item; an executor claims it with \texttt{workflowWorkItem.claim}. \\ \hline
\texttt{event} & Reserved for EventHandler notification; currently handled as \texttt{pull}. \\ \hline
\end{longtable}
}

\msubsection{enum}{Rule}
Transition rules: the eight rules of~\cite{varga2026} and the extensions (\texttt{x-}).

{\small
\begin{longtable}{|p{3.9cm}|p{11.9cm}|}

\hline
\rowcolor{gray!33} Value & Meaning \\ \hline
\endhead
\texttt{A-Start} & An action becomes active after its predecessor completed. \\ \hline
\texttt{S-Enable} & A step becomes enabled when its action is active and its predecessors completed. \\ \hline
\texttt{S-Bind} & An enabled step is bound to a provider (and executor) and becomes running. \\ \hline
\texttt{S-Done} & A running step completed; its result is stored. \\ \hline
\texttt{A-Finish} & Every step of an action completed. \\ \hline
\texttt{Z-Close} & Every action completed: the session is completed. \\ \hline
\texttt{S-FailBind} & No executor/provider could be bound: the session fails. \\ \hline
\texttt{S-FailInvoke} & The invocation failed: the session fails. \\ \hline
\texttt{x-rebind} & Extension: the step is bound again after \texttt{ERROR} or an unreachable executor. \\ \hline
\texttt{x-abort} & Extension: the session was aborted. \\ \hline
\texttt{x-timeout} & Extension: a step exceeded its timeout. \\ \hline
\texttt{x-skip} & Extension: a step was skipped because no guarded transition to it was taken. \\ \hline
\end{longtable}
}

\msubsection{struct}{Session}
The state $z = (M, M_A, \beta, \varepsilon, \xi, status)$ of a session, as seen from outside.

{\small
\begin{longtable}{|p{3.9cm}|p{4.4cm}|p{7.5cm}|}

\hline
\rowcolor{gray!33} Field & Type & Description \\ \hline
\endhead
\texttt{id} & \pref{SessionId} & Session identifier. \\ \hline
\texttt{recipe} & \pref{Name} & Recipe name. \\ \hline
\texttt{version} & \pref{Version} & Recipe version. \\ \hline
\texttt{subject} & \pref{String} & Subject of the session. Optional. \\ \hline
\texttt{owner} & \pref{SystemName} & The system that started the session; it may abort it. \\ \hline
\texttt{status} & \mref{SessionStatus} & \texttt{running}, \texttt{completed} or \texttt{failed}. \\ \hline
\texttt{dispatch} & \mref{DispatchMode} & Effective dispatch mode. \\ \hline
\texttt{failure} & \pref{String} & Why the session failed. Optional. \\ \hline
\texttt{extensions} & \pref{List}$<$\pref{String}$>$ & Extensions used (\texttt{x-rebind}, \texttt{x-abort}, \texttt{x-timeout}, \texttt{x-skip}, \ldots); empty when the run stays inside the proven model. \\ \hline
\texttt{stepMarking} & \pref{Map}$<$\pref{Name}, \pref{String}$>$ & $M$: \texttt{inactive}, \texttt{enabled}, \texttt{running}, \texttt{completed} or \texttt{skipped} per step. \\ \hline
\texttt{actionMarking} & \pref{Map}$<$\pref{Name}, \pref{String}$>$ & $M_A$: \texttt{inactive}, \texttt{active} or \texttt{completed} per action. \\ \hline
\texttt{modelTransitions} & \pref{Integer} & Number of model transitions so far. \\ \hline
\texttt{transitionBound} & \pref{Integer} & Theorem~1 bound $3|S| + 2|A|$. \\ \hline
\texttt{product} & \pref{SystemName} & Product system, when the session came from one. Optional. \\ \hline
\texttt{createdAt} & \pref{DateTime} & Start time. \\ \hline
\texttt{updatedAt} & \pref{DateTime} & Time of the last transition. \\ \hline
\end{longtable}
}

\msubsection{struct}{SessionFilter}
Optional filter of \oref{query}.

{\small
\begin{longtable}{|p{3.9cm}|p{4.4cm}|p{7.5cm}|}

\hline
\rowcolor{gray!33} Field & Type & Description \\ \hline
\endhead
\texttt{status} & \mref{SessionStatus} & Only sessions in this status. Optional. \\ \hline
\end{longtable}
}

\msubsection{struct}{SessionList}
Sessions, newest first.

{\small
\begin{longtable}{|p{3.9cm}|p{4.4cm}|p{7.5cm}|}

\hline
\rowcolor{gray!33} Field & Type & Description \\ \hline
\endhead
\texttt{entries} & \pref{List}$<$\mref{Session}$>$ & The sessions. \\ \hline
\texttt{count} & \pref{Integer} & Number of entries. \\ \hline
\end{longtable}
}

\msubsection{struct}{SessionRef}
Identifies a session; \oref{abort} also takes an optional reason.

{\small
\begin{longtable}{|p{3.9cm}|p{4.4cm}|p{7.5cm}|}

\hline
\rowcolor{gray!33} Field & Type & Description \\ \hline
\endhead
\texttt{id} & \pref{SessionId} & Session identifier. \\ \hline
\texttt{reason} & \pref{String} & Reason recorded with an abort. Optional. \\ \hline
\end{longtable}
}

\msubsection{enum}{SessionStatus}
Status of a session. \texttt{completed} and \texttt{failed} are final.

{\small
\begin{longtable}{|p{3.9cm}|p{11.9cm}|}

\hline
\rowcolor{gray!33} Value & Meaning \\ \hline
\endhead
\texttt{running} & Some rule may still apply. \\ \hline
\texttt{completed} & Z-Close fired: every action completed. \\ \hline
\texttt{failed} & S-FailBind, S-FailInvoke or an extension (abort, timeout) ended the session. \\ \hline
\end{longtable}
}

\msubsection{struct}{StartSessionRequest}
Starts one session (one execution of one recipe version for one subject).

{\small
\begin{longtable}{|p{3.9cm}|p{4.4cm}|p{7.5cm}|}

\hline
\rowcolor{gray!33} Field & Type & Description \\ \hline
\endhead
\texttt{recipe} & \pref{Name} & Recipe name; may be omitted when \texttt{product} is given. Optional. \\ \hline
\texttt{version} & \pref{Version} & Recipe version; the latest active version when absent. Optional. \\ \hline
\texttt{subject} & \pref{String} & What the session works on: product unit, control period or batch. Optional. \\ \hline
\texttt{parameters} & \pref{JsonObject} & Session parameters $\rho$; required parameters must be present. Optional. \\ \hline
\texttt{dispatch} & \mref{DispatchMode} & Overrides the recipe's dispatch mode. Optional. \\ \hline
\texttt{product} & \pref{SystemName} & A product system providing \texttt{workflowProduct}: the recipe and parameters are retrieved from it, and the results are stored back at Z-Close. Optional. \\ \hline
\end{longtable}
}

\msubsection{struct}{Trace}
The run of a session: the transitions and the data the rules produced.

{\small
\begin{longtable}{|p{3.9cm}|p{4.4cm}|p{7.5cm}|}

\hline
\rowcolor{gray!33} Field & Type & Description \\ \hline
\endhead
\texttt{sessionId} & \pref{SessionId} & Session identifier. \\ \hline
\texttt{status} & \mref{SessionStatus} & Current status. \\ \hline
\texttt{trace} & \pref{List}$<$\mref{Transition}$>$ & Transitions in order. \\ \hline
\texttt{results} & \pref{Map}$<$\pref{Name}, \pref{JsonValue}$>$ & $\xi$: stored step results. \\ \hline
\texttt{bindings} & \pref{Map}$<$\pref{Name}, \pref{JsonValue}$>$ & $\beta$: orchestration result per bound step. \\ \hline
\texttt{executorAssignment} & \pref{Map}$<$\pref{Name}, \pref{SystemName}$>$ & $\varepsilon$: executor per step. \\ \hline
\texttt{runner} & \pref{Map}$<$\pref{Name}, \pref{SystemName}$>$ & The system that ran each step. \\ \hline
\texttt{modelTransitions} & \pref{Integer} & Transitions of the proven rule set. \\ \hline
\texttt{transitionBound} & \pref{Integer} & Theorem~1 bound. \\ \hline
\texttt{withinModel} & \pref{Boolean} & True when no extension was used, so Theorems~1--2 apply. \\ \hline
\end{longtable}
}

\msubsection{struct}{Transition}
One transition of the session.

{\small
\begin{longtable}{|p{3.9cm}|p{4.4cm}|p{7.5cm}|}

\hline
\rowcolor{gray!33} Field & Type & Description \\ \hline
\endhead
\texttt{seq} & \pref{Integer} & Sequence number from 1. \\ \hline
\texttt{rule} & \mref{Rule} & The rule that fired. \\ \hline
\texttt{target} & \pref{Name} & The step or action it applied to. Optional. \\ \hline
\texttt{detail} & \pref{String} & Additional information, e.g. the executor or failure. Optional. \\ \hline
\texttt{at} & \pref{DateTime} & When it fired. \\ \hline
\end{longtable}
}


\subsection{Primitives}
\label{sec:model:primitives}
Types and structures mentioned throughout this document that are assumed to be available to implementations of this
service. The concrete interpretation of each of these types must be provided by any IDD document claiming to implement
this service.

{\small
\begin{longtable}{|p{3.2cm}|p{12.6cm}|}

\hline
\rowcolor{gray!33} Type & Description \\ \hline
\endhead
\pdef{Boolean} & One out of \texttt{true} or \texttt{false}. \\ \hline
\pdef{DateTime} & Pinpoints a moment in time (UTC). \\ \hline
\pdef{Integer} & A whole number. \\ \hline
\pdef{JsonObject} & A set of named values. \\ \hline
\pdef{JsonValue} & Any value: object, array, string, number, Boolean or null. \\ \hline
\pdef{List} & An ordered collection of items of the type given in angle brackets. \\ \hline
\pdef{Map} & A collection of key/value pairs of the types given in angle brackets. \\ \hline
\pdef{Name} & An identifier of a recipe, action, step, parameter or resource: a letter followed by at most 62 letters, digits, \texttt{-} or \texttt{\_}. \\ \hline
\pdef{SessionId} & The identifier the WorkflowManager gives a session. \\ \hline
\pdef{String} & A chain of characters. \\ \hline
\pdef{SystemName} & The name of an Arrowhead system as it appears in its certificate and in the ServiceRegistry. \\ \hline
\pdef{Version} & A recipe version in the form MAJOR.MINOR.PATCH. \\ \hline
\end{longtable}
}


\newpage
\bibliographystyle{IEEEtran}
\bibliography{bibliography}

\newpage

\section{Revision History}
\subsection{Amendments}

\noindent\begin{tabularx}{\textwidth}{| p{1cm} | p{2.2cm} | p{1.6cm} | X | p{3.6cm} |} \hline
\rowcolor{gray!33} No. & Date & Version & Subject of Amendments & Author \\ \hline
1 & 2026-11-23 & \arrowversion & Initial edition: service description of workflowInstance, written from the Java reference implementation (ah5-Workflow-java). & Documentation maintenance \\ \hline
\end{tabularx}

\subsection{Quality Assurance}

\noindent\begin{tabularx}{\textwidth}{| p{1cm} | p{2.2cm} | p{1.6cm} | X |} \hline
\rowcolor{gray!33} No. & Date & Version & Approved by \\ \hline
1 & & \arrowversion & Pending review \\ \hline
\end{tabularx}

\end{document}
